% Figure area_law_buffer_contours.
% Owner: chapter/ch24_peps_area_law_initial_buffer.tex (thm:area_law_buffer_contours).
% Used by: gapped-2d-notes AL Sec. 4 (Lemma 3.2 of the area-law paper, Move 1).
\begin{tenkzequation}
    % Source: arXiv OpenAI2026AreaLaw, sections/02-initial.tex:219--300
    % (lemma lem:initial-buffer and the contour separation in its proof);
    % the source has no figure.
    % Geometry only: dots are sites of Q (nearest-neighbor edges are not
    % drawn); there are no tensors or open indices, and the boundary
    % signature is empty.
    % Enclosures, inner to outer: Q_0 = rows/cols 4..5; Q_0^{+d} with
    % d=1 = rows/cols 3..6; Q_0^{+d'} with d'=2 = rows/cols 2..7; the
    % safe rectangle Q = rows/cols 1..8. All of Q is drawn inside A.
    % Tinted pairs: adjacent pairs, admissible at range R=1. The pair
    % (2,4),(3,4) is split by Q_0^{+d} and lies in Q_0^{+(d+R)} = Q_0^{+d'};
    % the pair (1,6),(2,6) is split by Q_0^{+d'} only. Since d+R<=d',
    % neither is split by both contours.
    % Not to scale: safety puts the cut endpoints Z at sup-norm distance
    % greater than D_0 size(Q) from Q, so Z lies far outside the picture.
    \tenkzkernel
    \begin{tenkz}[size=s, rows={wire,wire,wire,wire,wire,wire,wire,wire},
            cols=8, bonds=none]
        \tndeclare{species}{support}{hue=source:orange!85!black}
        \tn{} & \tn{} & \tn{} & \tn{} & \tn{} & \tn{} & \tn{} & \tn{} \\
        \tn{} & \tn{} & \tn{} & \tn{} & \tn{} & \tn{} & \tn{} & \tn{} \\
        \tn{} & \tn{} & \tn{} & \tn{} & \tn{} & \tn{} & \tn{} & \tn{} \\
        \tn{} & \tn{} & \tn{} & \tn{} & \tn{} & \tn{} & \tn{} & \tn{} \\
        \tn{} & \tn{} & \tn{} & \tn{} & \tn{} & \tn{} & \tn{} & \tn{} \\
        \tn{} & \tn{} & \tn{} & \tn{} & \tn{} & \tn{} & \tn{} & \tn{} \\
        \tn{} & \tn{} & \tn{} & \tn{} & \tn{} & \tn{} & \tn{} & \tn{} \\
        \tn{} & \tn{} & \tn{} & \tn{} & \tn{} & \tn{} & \tn{} & \tn{}
        \tnmark[form=enclosure, species=selected, tint, label pos=e]{(4,4) .. (5,5)}{$Q_0$}
        \tnmark[form=enclosure, species=secondary, label pos=e]{(3,3) .. (6,6)}{$Q_0^{+d}$}
        \tnmark[form=enclosure, species=complement, label pos=e]{(2,2) .. (7,7)}{$Q_0^{+d'}$}
        \tnmark[form=enclosure, label pos=e]{(1,1) .. (8,8)}{$Q$}
        \tnmark[form=enclosure, species=support, tint]{{(2,4), (3,4)}}{}
        \tnmark[form=enclosure, species=support, tint]{{(1,6), (2,6)}}{}
    \end{tenkz}
\end{tenkzequation}
